\chapter{典型应用场景与用户需求}

\section{场景一：工业设备与生产服务监测}

\subsection{现场环境与核心任务}

在车间、厂区或边缘站点中，设备控制程序、工控机操作系统、数据采集服务和生产应用共同形成一条现场业务链。它们持续输出设备状态、通信异常、任务执行、资源占用和服务错误等日志。日志来源彼此异构，但运维人员真正关心的不是某一行文本属于哪个类别，而是：异常是否正在影响生产、能否立即采取动作、是否需要调动更强的分析能力。

这类场景具有明确的时间与可用性约束。温度越限、磁盘写入失败、进程退出、设备连接中断等已知风险，应在日志抵达边缘节点后立即形成现场判断，不能等待公网往返或云端模型排队；设备编号、路径、地址和业务参数等内容又可能受到数据出域策略限制。与此同时，版本升级、工艺切换和设备老化会不断产生新模板，单纯依赖预置规则难以覆盖长期运行中的开放异常。因此，系统既要保留规则路径的确定性和低时延，又要为未知事件提供逐级增强的语义分析能力。

\subsection{典型故障链与系统处理过程}

以“设备数据采集服务异常”为例，故障可能依次表现为驱动连接超时、采集任务重试、本地队列积压和上层生产服务响应变慢。若逐条查看，多个日志看似来自不同进程；若全部直接上云，又会把大量正常心跳和重复重试一并传输。CoTriage 对这一过程按事件价值分层处理：采集器先把日志转换为带有设备、来源、时间和稳定标识的统一事件；规则与模板优先识别明确故障，轻量分类器和边缘小语言模型处理未命中或语义更复杂的内容；RAISE 再依据严重度、不确定性、新颖度、跨设备影响与实时网络状态选择 \texttt{LOCAL}、\texttt{CLOUD} 或 \texttt{DEFER}。

若事件已被边缘知识充分解释，系统直接向现场输出异常类型、置信度和推理路径；若出现新模板、连续异常或多个服务同时受影响，云端围绕主事件聚合有限的时间邻近与同源上下文，生成根因候选、处置建议和解释；若网络不可用但事件仍值得深诊，现场初判照常产生，上云请求由 Outbox 持久化并等待恢复后按风险优先重放。由此，云端分析是现场判断的增益通道，而不是实时感知成立的前置条件。

\subsection{业务价值与可验收结果}

该场景最终向现场人员交付三类可操作信息：第一，正在发生什么以及影响对象；第二，系统为何在本地处理、立即上云或延期上云；第三，复杂事件可能的根因与建议动作。管理人员还可查看事件经历的规则、模板或模型版本，核对敏感数据是否被限制出域。因而，该场景的验收重点是边缘结果能否持续产生、已知告警能否快速收敛、困难事件能否获得增量解释，以及断网期间高价值事件是否仍被可靠保留，而不是只比较单个分类器的离线分数。

\section{场景二：分布式基础设施监测}

\subsection{运行环境与诊断难点}

在 Hadoop、ZooKeeper、Linux 服务与 Web 服务组成的分布式基础设施中，一个业务请求通常跨越多个进程和节点。组件异步交互、失败重试和副本机制会使根因与症状分离：依赖服务短暂抖动，可能在上游表现为连接超时，在计算节点表现为任务失败，在协调节点表现为会话重建。运维人员面对的不是一条决定性日志，而是来自不同时间、不同节点、不同严重度的碎片化证据。

传统单节点告警容易把传播症状分别判成多个独立故障，形成告警风暴；全量集中后再分析虽然获得更多信息，却增加传输和检索负担，而且大量重复日志会稀释关键证据。该场景因此更强调事件关联、上下文裁剪和判断可解释性：系统既要保留每个节点的快速初判，又要在必要时恢复跨节点的时间关系和依赖关系，帮助区分单点组件故障、网络分区、资源争用、共同依赖异常和集群级事件。

\subsection{典型故障链与系统处理过程}

以“ZooKeeper 会话异常引发上游服务抖动”为例，协调节点可能先出现连接或选举相关事件，随后多个应用节点产生会话过期、重连和请求超时。各边缘节点独立采集并完成初判，统一事件结构保留设备、来源、模板、严重度和时间坐标；重复且已知的重连信息可留在本地，高严重度、新颖或多节点共同出现的事件则获得更高云端增益分数。云端网关以选中的主事件为中心，按时间邻近、同源、模板与跨设备关系裁剪上下文，再请求大模型判断故障起点、传播方向和优先处置对象。

诊断输出不是简单复述日志，而是把“现象—证据—根因候选—建议动作”组织为结构化结果。例如，系统可以说明多个应用节点的超时集中发生在协调服务异常之后，因此应优先检查协调节点与网络连通性，而不是逐台重启上游应用。云端结论写回后与原事件、路由快照和上下文关联，运维人员能够回看诊断采用了哪些证据；若边云判断存在差异，相关统计还可作为后续路由决策的反馈信号。

\subsection{扩展价值与可验收结果}

该场景验证的是系统从“单点日志分类”走向“分布式事件分诊”的能力。新增节点时，节点仍使用相同的事件契约接入，路由控制器依据事件状态作出决策，云端按需聚合相关上下文；节点数量和日志量增长时，采集、推理、路由与模型服务可以独立部署和扩容。验收时可沿一个跨节点故障观察：多个来源是否被统一登记，边缘是否保留各自初判，关键事件是否被正确路由，云端解释是否包含关联证据，以及网络异常时延期请求是否进入可追踪的 Outbox 状态。

\section{两类场景的共性主链与差异化适配}

工业现场与分布式基础设施在设备类型、故障传播方式和处置责任上差异明显，但共享“连续事件流、局部快速判断、少量复杂事件深诊、弱网下保持主链”这一业务抽象。CoTriage 将稳定能力下沉为统一事件、分层推理契约、RAISE、Outbox、云端网关和管理接口，将场景差异限制在采集适配、领域规则、模板知识、模型配置与上下文关联策略中。

\begin{table}[htbp]
  \centering
  \small
  \caption{两类典型场景的差异化需求与共用能力}
  \label{tab:scenario-comparison}
  \begin{tabularx}{\textwidth}{>{\bfseries}p{0.16\textwidth} X X}
    \toprule
    对比维度 & 工业设备与生产服务 & 分布式基础设施 \\
    \midrule
    主要对象 & 设备程序、工控机、采集服务与生产应用 & Hadoop、ZooKeeper、Linux 与 Web 服务节点 \\
    典型风险 & 设备离线、采集失败、进程异常与生产任务受阻 & 网络分区、资源争用、依赖异常与集群级故障 \\
    首要约束 & 现场低时延、断网可用、数据边界与处置明确 & 跨节点关联、传播分析、诊断解释与规模扩展 \\
    场景适配项 & 设备字段、领域规则、模板知识与出域策略 & 节点关系、服务模板、上下文窗口与关联策略 \\
    共用主链 & \multicolumn{2}{p{0.75\textwidth}}{统一事件化、边缘分层推理、RAISE 三动作路由、风险优先 Outbox、云端结构化诊断与统一展示} \\
    \bottomrule
  \end{tabularx}
\end{table}

这种“场景适配局部化、协同主链稳定化”的方式使扩展不等于复制系统：新增日志类型时无需重写路由、可靠补偿、云端调用和展示逻辑；更换边缘模型或云端 Provider 时也不改变上游事件契约。它既降低第二类场景接入成本，也为后续向能源站点、园区基础设施和其他分布式运维场景迁移保留了稳定边界。

\section{角色与典型任务}
\begin{table}[htbp]
  \centering
  \small
  \caption{系统角色、典型任务与主要关注点}
  \label{tab:roles-v03}
  \begin{tabularx}{\textwidth}{>{\bfseries}p{0.18\textwidth} X X}
    \toprule
    角色 & 典型任务 & 主要关注点 \\
    \midrule
    现场运维人员 & 查看实时事件、异常态势、路由解释和诊断建议 & 响应速度、易读性、断网可用 \\
    系统管理员 & 管理设备、规则、模板、边缘模型和云端 Provider & 配置版本、启停状态、权限和回滚 \\
    平台开发人员 & 接入新日志来源、模型服务与下游处置系统 & 接口稳定、幂等、可追踪和扩展成本 \\
    评审与验收人员 & 沿演示路径验证采集、推理、路由、深诊和展示 & 功能完整、解释清晰、口径一致 \\
    \bottomrule
  \end{tabularx}
\end{table}

\section{功能需求}

功能需求沿一条事件从产生到闭环的生命周期展开，而不是按页面菜单罗列。系统首先要把异构日志转化为稳定事件，随后形成边缘初判和路由动作；只有被选中的事件才进入云端深诊，暂不可发送的事件必须留在可靠补偿链中；整个过程还要能够被配置、查询和解释。\cref{tab:functional-requirements}给出了可由后续组件和页面逐项承接的功能需求。

\begin{table}[htbp]
  \centering
  \footnotesize
  \caption{CoTriage 核心功能需求与可观察结果}
  \label{tab:functional-requirements}
  \begin{tabularx}{\textwidth}{>{\bfseries\ttfamily}p{0.10\textwidth} >{\bfseries}p{0.21\textwidth} X p{0.24\textwidth}}
    \toprule
    \normalfont 编号 & 能力 & 系统行为 & 可观察结果 \\
    \midrule
    FR-01 & 多节点采集与事件化 & 持续读取节点新增日志，生成稳定事件标识并补充设备、来源、时间和原文 & 设备可登记；事件可按来源检索并进入原始事件流 \\
    FR-02 & 边缘分层推理 & 按规则、模板、轻量分类器和边缘小语言模型由快到慢处理，保留类别、严重度、置信度、证据与版本 & 事件详情能够显示实际命中层级、推理路径和处理时延 \\
    FR-03 & 自适应协同路由 & RAISE 联合风险、不确定性、新颖度、网络、队列、预算和策略，输出 \texttt{LOCAL/CLOUD/DEFER} & 每次决策保留动作、分数、动态门槛、原因和状态快照 \\
    FR-04 & 弱网可靠补偿 & 将暂不可上云但仍有价值的事件写入 Outbox，按风险优先级和到期时间重试，超限后进入死信 & 可查询待处理、重试、完成和死信状态，并追踪最近错误 \\
    FR-05 & 云端关联诊断 & 围绕主事件筛选相关上下文，脱敏并控制 Token 预算，生成根因候选、处置建议和结构化解释 & 诊断结果能够关联回原事件，并展示 Provider、耗时和上下文依据 \\
    FR-06 & 反馈与状态写回 & 将云端纠错、信息增量和高风险确认写回事件状态与路由统计，避免诊断成为孤立文本 & 边缘结果、路由记录、云端诊断和反馈可按同一事件标识串联 \\
    FR-07 & 在线配置治理 & 提供设备、规则、模板、轻量模型、边缘小模型和云端 Provider 的查询、校验、发布与状态管理 & 管理端显示版本、生效状态和遮罩后的敏感配置，失败发布不替换有效版本 \\
    FR-08 & 统一观测与下钻 & 通过 REST/SSE 提供历史快照和实时增量，集中展示事件、异常、路由、Outbox、模型与组件状态 & 用户可从总览下钻至单事件处理链，并区分运行指标与实验结论 \\
    \bottomrule
  \end{tabularx}
\end{table}

这些需求共同构成四个闭环：感知闭环将日志转为事件，决策闭环将边缘结果转为三动作路由，诊断闭环将高价值上下文转为结构化结论，治理闭环负责配置、追踪和解释。任何单项能力都不能脱离统一事件标识独立存在，否则系统会退化为彼此割裂的模型与页面演示。

\section{非功能需求}

非功能需求决定系统在压力、故障和配置变化下是否仍然可信。由于本作品同时包含实时数据面、动态决策面和在线控制面，不能只以“接口成功返回”作为质量标准；还必须规定关键依赖失效时哪些能力继续运行、哪些任务可以延期、哪些数据不得越过边界，以及事后能否还原当时的系统状态。

\begin{table}[htbp]
  \centering
  \footnotesize
  \caption{CoTriage 非功能需求、约束与退化策略}
  \label{tab:nonfunctional-requirements}
  \begin{tabularx}{\textwidth}{>{\bfseries\ttfamily}p{0.10\textwidth} >{\bfseries}p{0.17\textwidth} X X}
    \toprule
    \normalfont 编号 & 质量属性 & 核心约束 & 异常时的退化与验证方式 \\
    \midrule
    NFR-01 & 实时性 & 边缘快速路径不得等待云端完成；采集、推理、路由与云端阶段分别计时 & 云端变慢时仍先产生边缘结果；通过事件耗时与窗口 P95 观察尾时延 \\
    NFR-02 & 可用性与弱网韧性 & 网络或 Provider 不可用不得阻塞现场感知，高价值延期任务不得静默丢失 & 保留 \texttt{LOCAL} 结果，将待上云事件转为 \texttt{DEFER} 并持久化 \\
    NFR-03 & 可靠性与幂等 & 当前阶段持久化及下游发布完成后才能确认消息；相同幂等键不得重复触发高成本调用 & 重试复用既有状态；失败进入 Pending、Outbox 或 DLQ，并保留次数与错误 \\
    NFR-04 & 一致性与可追溯 & 原文、推理、路由、诊断和反馈必须通过稳定事件标识及版本串联 & 能够还原当次模型、规则、门槛和系统快照，区分旧结果与配置更新后的结果 \\
    NFR-05 & 安全与隐私 & 上云数据遵循最小化与脱敏原则；明确的出域禁令和生产策略优先于模型评分 & 命中硬策略时禁止上云；密钥只展示遮罩值，诊断上下文保留脱敏边界 \\
    NFR-06 & 可扩展性 & 采集、推理、路由、网关、模型和管理端通过稳定事件或 HTTP 契约解耦 & 节点或吞吐增长时独立增加对应实例；接入新来源不重写协同主链 \\
    NFR-07 & 配置安全性 & 配置发布必须完整解析和校验，运行组件只原子切换到完整有效版本 & 拉取或校验失败时继续使用上一版本，不加载半份配置并暴露失败状态 \\
    NFR-08 & 可观测与可维护 & 健康、Stream、Pending、队列、预算、Outbox、DLQ 和 Provider 状态必须可查询 & 运维人员能够由看板定位到组件和单事件，页面不可用不影响后台数据面 \\
    \bottomrule
  \end{tabularx}
\end{table}

上述约束定义了 CoTriage 的“受控退化”原则：云端不可用时退化分析深度而不停止边缘判断，模型不可用时退回已有知识层而不伪造推理结果，配置失败时保留上一有效版本而不切换到不完整状态，消息处理失败时进入可追踪的重试或死信边界而不静默丢弃。后续章节将分别把这些要求落实到统一事件、消息拓扑、分层推理、RAISE、Outbox、云端网关和管理中心中。
